\subsection{Droits et vues}

\subsubsection{Profils et convention d'accès}

Trois profils utilisent la base. L'administrateur gère l'ensemble des données. L'organisateur assure le suivi des participants, des équipes, des inscriptions, du calendrier et des résultats. Le public consulte uniquement les informations publiées.

Les tableaux distinguent l'accès direct aux données et leur consultation à travers des vues. L signifie lecture, I insertion, M modification et S suppression. « Via vue » autorise uniquement la lecture des informations décrites dans les vues ci-dessous ; le public ne dispose d'aucun accès direct aux données sous-jacentes. Les vues sont consultables par l'organisateur et l'administrateur, ainsi que par le public lorsque cela est indiqué.

Les droits de gestion restent soumis à toutes les contraintes métier. En particulier, les inscriptions sont conservées pendant les Jeux, leur date initiale et le couple équipe--épreuve sont immuables, et les équipes et épreuves correspondantes ne peuvent pas être supprimées. La composition d'une équipe est figée dès sa première inscription. Ces règles s'appliquent aussi à l'administrateur.

\subsubsection{Droits sur les concepts}

\begin{table}[H]
  \centering
  \footnotesize
  \renewcommand{\arraystretch}{1.3}
  \begin{tabular}{|>{\raggedright\arraybackslash}p{3.1cm}|
                  >{\centering\arraybackslash}p{3.0cm}|
                  >{\centering\arraybackslash}p{4.4cm}|
                  >{\centering\arraybackslash}p{3.0cm}|}
    \hline
    \rowcolor{gray!25}
    \textbf{Concept} & \textbf{Admin.} & \textbf{Organisateur} & \textbf{Public} \\
    \hline
    \textbf{Athlète} & L, I, M, S & L, I, M, S & Via vue \\
    \hline
    \textbf{Équipe} & L, I, M, S & L, I, M, S & Via vue \\
    \hline
    \textbf{Entraîneur} & L, I, M, S & L, I, M, S & Via vue \\
    \hline
    \textbf{Discipline} & L, I, M, S & L, I, M, S & Via vue \\
    \hline
    \textbf{Épreuve} & L, I, M, S & L, I, M, S & Via vue \\
    \hline
    \textbf{Lieu} & L, I, M, S & L, I, M, S & Via vue \\
    \hline
  \end{tabular}
  \caption{Droits par concept, sous réserve des contraintes métier}
  \label{tab:droits-concepts}
\end{table}

Le public ne voit ni le poids, ni la taille, ni la date de naissance des athlètes. Pour les entraîneurs, seuls le nom, le prénom, la nationalité et les spécialités sont publiés ; la date de naissance et le numéro de carte professionnelle restent réservés à la gestion.

\clearpage
\subsubsection{Droits sur les associations}

L'organisateur et l'administrateur peuvent constituer les équipes, renseigner les disciplines pratiquées et les spécialités, puis créer, annuler ou réactiver les inscriptions avant le début des épreuves. L'annulation définitive d'une épreuve entraîne automatiquement celle de ses inscriptions actives, sans supprimer leur historique.

{\footnotesize
\renewcommand{\arraystretch}{1.3}
\begin{longtable}{|>{\raggedright\arraybackslash}p{3.1cm}|
                  >{\centering\arraybackslash}p{3.0cm}|
                  >{\centering\arraybackslash}p{4.4cm}|
                  >{\centering\arraybackslash}p{3.0cm}|}
  \caption{Droits par association, sous réserve des contraintes métier}
  \label{tab:droits-associations} \\
  \hline
  \rowcolor{gray!25}
  \textbf{Association} & \textbf{Admin.} & \textbf{Organisateur} & \textbf{Public} \\
  \hline
  \endfirsthead
  \hline
  \rowcolor{gray!25}
  \textbf{Association} & \textbf{Admin.} & \textbf{Organisateur} & \textbf{Public} \\
  \hline
  \endhead
  \hline
  \endfoot
  \textbf{Composé}\newline (athlète -- équipe) & L, I, M, S & L, I, M, S & Aucun \\
  \hline
  \textbf{Encadre}\newline (entraîneur -- équipe) & L, I, M, S & L, I, M, S & Aucun \\
  \hline
  \textbf{Spécialisé}\newline (entraîneur -- discipline) & L, I, M, S & L, I, M, S & Via vue \\
  \hline
  \textbf{Pratique}\newline (athlète -- discipline) & L, I, M, S & L, I, M, S & Via vue \\
  \hline
  \textbf{Concourt}\newline (équipe -- discipline) & L, I, M, S & L, I, M, S & Via vue \\
  \hline
  \textbf{Participe}\newline (équipe -- épreuve) & L, I, M du statut & L, I, M du statut & Aucun \\
  \hline
  \textbf{Comporte}\newline (épreuve -- discipline) & L, I, M, S & L, I, M, S & Via vue \\
  \hline
  \textbf{Accueille}\newline (lieu -- épreuve) & L, I, M, S & L, I, M, S & Via vue \\
  \hline
  \textbf{Résultat}\newline (équipe -- épreuve) & L, I, M ; S des provisoires & L, I, M ; S des provisoires & Via vue, validés uniquement \\
  \hline
\end{longtable}
}

Pour un résultat, l'insertion crée toujours un état provisoire. La modification des données du résultat est autorisée uniquement à cet état. La validation et le retour au statut provisoire sont réservés à l'organisateur et à l'administrateur ; un résultat validé ne peut pas être supprimé directement. Son retour au statut provisoire le retire d'abord de la publication et du tableau des médailles.

\clearpage
\subsubsection{Vues}

Les vues publiques présentent uniquement les attributs indiqués ci-dessous. Elles relient les informations utiles ou calculent des totaux, sans donner accès aux autres données de gestion.

{\footnotesize
\renewcommand{\arraystretch}{1.3}
\begin{longtable}{|>{\raggedright\arraybackslash}p{4.1cm}|
                  >{\raggedright\arraybackslash}p{7.0cm}|
                  >{\centering\arraybackslash}p{3.0cm}|}
  \caption{Vues de la base de données}
  \label{tab:vues} \\
  \hline
  \rowcolor{gray!25}
  \textbf{Vue} & \textbf{Contenu} & \textbf{Profils concernés} \\
  \hline
  \endfirsthead
  \hline
  \rowcolor{gray!25}
  \textbf{Vue} & \textbf{Contenu} & \textbf{Profils concernés} \\
  \hline
  \endhead
  \hline
  \endfoot
  \texttt{\scriptsize V\_ATHLETE\_PUBLIC} & Identifiant, nom, prénom, sexe et nationalité de l'athlète, accompagnés des noms des disciplines pratiquées. Une ligne par discipline ; les athlètes sans discipline restent visibles. Ni date de naissance, ni poids, ni taille. & Tous \\
  \hline
  \texttt{\scriptsize V\_ENTRAINEUR\_PUBLIC} & Nom, prénom et nationalité de l'entraîneur, accompagnés des noms de ses disciplines de spécialisation. Une ligne par spécialité ; les entraîneurs sans spécialité restent visibles. Ni date de naissance, ni numéro de carte professionnelle. & Tous \\
  \hline
  \texttt{\scriptsize V\_CALENDRIER} & Identifiant de l'épreuve, discipline, phase, sexe concerné, attribution de médaille, date, horaires, nom et adresse du lieu, pour les épreuves non annulées. & Tous \\
  \hline
  \texttt{\scriptsize V\_RESULTATS\_VALIDES} & Identifiant de l'équipe, discipline, épreuve (identifiant, phase et date), classement éventuel, issue sportive et médaille éventuelle. Uniquement les résultats validés. & Tous \\
  \hline
  \texttt{\scriptsize V\_TABLEAU\_MEDAILLES} & Une ligne par équipe avec son identifiant, sa discipline et le nombre de médailles d'or, d'argent et de bronze, ainsi que leur total. Seuls les résultats validés comptent. Les équipes sans médaille figurent avec des totaux nuls ; une médaille collective compte une seule fois. & Tous \\
  \hline
  \texttt{\scriptsize V\_SUIVI\_PREPARATION} & Épreuves non annulées à venir sans inscription active, avec leur identifiant, leur discipline, leur date, leurs horaires et leur lieu. Les inscriptions annulées ne comptent pas comme des participants. & Organisateur, administrateur \\
  \hline
\end{longtable}
}
